\begin{abstract}
Logical reasoning systems need different kinds of evidence for true, false, and unknown queries, but most language-model pipelines still emphasize successful proofs and under-specify refutation and abstention. We introduce \model, a supervision format that trains a model to choose among three verifiable reasoning modes: \prove{} for proving the query, \refute{} for proving the opposite literal, and \abstain{} for emitting a failed-rule or missing-support witness. We pair this format with a symbolic verifier that checks whether generated traces are consistent with the theory closure. On a fixed three-seed 7B suite with only 4,096 training questions, \model{} does not maximize raw accuracy, but it consistently improves verifier-backed reasoning: joint label-plus-evidence accuracy rises from 39.9\% to 75.0\% in-domain and from 26.5\% to 55.0\% on depth-5, while raw accuracy remains close to proof-only supervision and within 1.5--3.6 points of answer-only. The mechanism is faithful abstention rather than mere abstention frequency: on gold-unknown queries, proof-only and \model{} predict \texttt{Unknown} at similar rates, but \model{} produces 1,444 faithful unknown explanations in-domain versus 0 for proof-only supervision. As training scale increases, the raw-accuracy gap largely disappears. At full train, \model{} reaches 98.2\% in-domain accuracy with 96.8\% joint accuracy, and on depth-5 it attains 93.8\% accuracy and 86.2\% joint accuracy on the fixed subset, slightly surpassing answer-only in raw accuracy. These results support a precise claim: countermodel-aware supervision is primarily a mechanism for faithful reasoning, and at sufficient scale it can deliver that benefit without sacrificing strong answer accuracy.
\end{abstract}
